\begin{claim}{B1.2}{Proved}{All admissible twists have the same Hodge numbers}
For every family in B1.1 whose cusp is the one-component hexagonal filling,
the direct images of the absolute cotangent sheaves are
\[
\begin{array}{c|ccc}
 &R^0f'_*&R^1f'_*&R^2f'_*\\\hline
\Omega^1_{X'}&\OO(-2)&\OO\oplus\OO(-1)^2&\OO\oplus\OO(-1)\\
\Omega^2_{X'}&\OO(-2)\oplus\OO(-1)&\OO(-2)\oplus\OO(-1)^2&\OO
\end{array}
\]
and $R^qf'_*\OO_{X'}=(\OO,\OO\oplus\OO(-1),\OO(-1))$.
In particular every admissible $(3,4)$ member with
$|12l_0-4l_1-3l_2|=1$ has the Hodge diamond of Part A.
\end{claim}
\begin{proof}
B1.1 identifies the logarithmic direct images and their connecting maps.
In Part A the cokernel of $\kappa_0$ is
$\C_{p_1}\oplus\C_{p_2}$, and $\kappa_1$ is surjective with kernel
$\OO\oplus\OO(-1)^2$. These statements therefore hold for $X'$.
The residue sheaf of its normalized divisor has higher direct images
\[
R^0f'_*\OO_{\widetilde D'}=\C_{p_0}\oplus\C_{p_1}\oplus\C_{p_2},
\quad R^1f'_*\OO_{\widetilde D'}=\C_{p_1}\oplus\C_{p_2},
\quad R^2=0.
\]
Indeed the finite reductions are bielliptic, and the cusp normalization
is rational. The base logarithmic residue is evaluation with weights
$(1,m_1,m_2)$; its kernel is $\OO(1)(-\Sigma)=\OO(-2)$.
On $H^1$ the residue of the invariant class $\gamma^{0,1}$ is
$m_j\gamma^{0,1}|_{S'_j}\ne0$. The other eigenline has zero residue.
The local rank-one image of $\kappa_0$ is its kernel, as follows either
from the residue diagram of Part A or from the eigenframe calculation
in B1.1. Thus the two finite skyscrapers cancel exactly. This proves
the first row, with no surviving torsion.

The proper flat map is Gorenstein of relative dimension two. Relative
duality for $F=\Omega^1_{X'}$, using
$F^*\otimes\omega_{X'/B}=\Omega^2_{X'}\otimes f'^*\OO(2)$,
gives the second row because all sheaves in the first row are locally
free. Coherent Leray on $\Pone$ gives every Hodge number. In particular
$h^{p,0}=0$ for $p>0$ follows from these direct images and duality,
so no separate invariance assumption about absolute forms is needed.
\end{proof}
